% !TeX root = ../main_chapter1_preview.tex
\chapter{绪论}
\label{chap:introduction}

随着柔性制造、机器人辅助手术和复杂装备维护等应用不断发展，机器人正在从封闭、重复、刚性程序执行单元，逐渐走向开放、交互和接触丰富的任务环境。在这类环境中，机器人不再只需要沿着预先设定的轨迹运动，而是需要在操作者意图、环境几何约束和接触力安全之间实时权衡。更具体地说，机器人既要能够接受人的监督修正，又不能让人的误触、震颤或不安全输入直接传递到末端；既要保持接触式作业所需的稳定力，又不能因力过大造成工件损伤、组织损伤或工具冲击。这样的控制问题具有明显的多目标、多约束和强交互特征，是本文研究的出发点。

本文关注的应用对象主要包括两类。第一类是机器人辅助微创操作，例如带远心运动约束的长工具操作、柔性组织接触、手术训练中的标准化操作任务等。该类任务要求工具轴线满足固定通道约束，同时在工具尖端与组织接触时保持安全交互力。第二类是通用工业约束接触任务，例如柔性密封条检测、复合材料壳体打磨、曲面抛光、橡胶件表面扫描、接触式边缘跟随和柔性工件去毛刺等。该类任务中，操作者通常具有丰富的工艺经验，能够根据视觉观察和加工效果对机器人轨迹进行在线修正；机器人则具有稳定执行、力反馈采集和长时间重复操作能力。两类场景虽然应用背景不同，但都可以抽象为“人监督、机器人执行、几何受限、接触力受约束”的人机协同控制问题。

从方法上看，本文以合作博弈和动态仲裁为主线。第三章面向刚性几何约束操作，研究人类监督意图与机器人自主目标之间如何分配权重；第四章面向柔性接触任务，研究位置跟踪目标与接触力调节目标之间如何动态协调；第五章进一步将二者融合，面向工业约束接触任务构造参考层人机仲裁与执行层力--位仲裁相结合的双仲裁控制框架。本章按照硕士论文绪论的写法，首先说明研究背景和实际意义，然后从人机共享控制、微分博弈与动态仲裁、约束接触与力--位协同、手术机器人和工业落地应用等方面综述相关研究，最后给出本文研究内容、创新点和全文组织结构。

\section{研究背景及意义}
\label{sec:ch1_background}

传统工业机器人多工作在围栏内，环境、工件和任务流程相对固定，控制系统主要追求重复定位精度和节拍效率。近年来，生产现场正在从大批量、标准化、刚性工装生产转向小批量、多品种、柔性化和人机协作生产。以复合材料风电叶片修复为例，打磨区域通常由缺陷位置决定，工件表面存在曲率变化和材料余量不均；以汽车密封条、橡胶件或软质覆盖件检测为例，机器人需要以较小接触力沿表面扫描，同时操作者需要根据局部外观和工艺经验修正轨迹；以曲面抛光和去毛刺为例，切向轨迹决定覆盖质量，法向接触力决定加工质量和工具寿命。RoboGrind 和增强现实示教编程等面向工业机器人程序生成的研究也表明，工业打磨、砂光和抛光任务具有自动化需求强、编程成本高、力控要求明显等特点，交互式编程和力控执行是其落地的重要环节\cite{alt2024robogrind,dogangun2025rampa}。

机器人辅助手术则从另一个方向提出了类似需求。微创手术中的长工具需要通过穿刺孔进入体内，工具轴线受到远心运动约束；末端与组织接触时，交互力过大可能造成组织损伤，过小又可能导致抓持、牵拉或接触感知失效。近年来，自主软组织手术、情境感知力控制和增强现实安全辅助等研究持续推进，说明机器人辅助手术正在从单纯遥操作走向更高等级的监督自主和安全辅助\cite{saeidi2022star,ishida2024force,ishida2024haptic,chen2023markerless}。不过，手术场景中仍然存在软组织变形复杂、临床数据不足、安全验证困难等问题，基础模型和多模态学习虽然提供了新的方向，但无法直接替代稳定、可解释和可分析的底层控制器\cite{schmidgall2024foundation,zargarzadeh2024bloodsuction}。

无论是工业柔性工件操作还是机器人辅助手术，控制问题都不再是单一目标。首先，操作者输入具有任务意义。人在视觉理解、异常判断、工艺经验和安全监督方面仍然具有优势。比如，当机器人沿复合材料表面扫描时，操作者可以发现局部缺陷区域并要求机器人改变扫描路径；当机器人进行接触式检测时，操作者可以根据传感器显示或视觉反馈调整重点区域。共享自主研究已经指出，机器人应当在理解人类意图的基础上提供辅助，而不是简单替代人类或完全服从瞬时输入\cite{dragan2013policy,jain2020probabilistic,jonnavittula2021learning,jonnavittula2022sari}。其次，接触力安全是刚性约束。人的输入可能沿切向修正路径，也可能不小心沿法向压入工件或组织。若控制器不区分二者，安全输入和危险输入都会被同等接受，容易造成力越界。

从力学角度看，柔性接触任务具有明显的力--位耦合。若将局部接触近似为弹簧--阻尼模型，接触力可写为
\begin{equation}
  F(t)=K_e\delta(t)+B_e\dot\delta(t)+\nu_f(t),
  \label{eq:ch1_kv_model}
\end{equation}
其中 $\delta$ 为压入深度，$\dot\delta$ 为压入速度，$K_e$ 和 $B_e$ 分别表示局部等效刚度和阻尼，$\nu_f$ 表示噪声、摩擦、接触角度变化和未建模柔性带来的残差。式\eqref{eq:ch1_kv_model} 的物理含义很直接：位置或速度的微小变化会被环境刚度和阻尼放大为接触力变化。若环境从低刚度区进入高刚度区，同样的位移误差会产生更大的力误差。因此，固定参数的位置控制或固定权重的力--位控制都难以在变刚度环境中稳定工作。

为了描述接触安全，通常需要设定接触力上下界
\begin{equation}
  F_{\min}<F_d<F_{\max},
  \label{eq:ch1_force_band}
\end{equation}
其中 $F_d$ 为期望接触力，$F_{\min}$ 为维持有效接触的下界，$F_{\max}$ 为安全上界。若 $F(t)<F_{\min}$，工具可能脱离工件或组织，导致扫描、检测或力控任务失败；若 $F(t)>F_{\max}$，则可能造成表面压痕、材料损伤、传感器过载或软组织风险。也就是说，接触控制不仅要减小力误差 $F-F_d$，还要识别当前力是否靠近上下边界。ForceMimic、FORGE 和学习式力控等近期研究都从不同角度说明，显式力信息对于接触丰富操作的鲁棒性和安全性具有重要价值\cite{beltran2020learning,liu2024forcemimic,noseworthy2024forge,portela2024learning}。

在约束操作中，还需要考虑工具、工装或通道引入的几何约束。设固定约束点为 $p_c$，工具轴线为 $\ell(t)$，几何约束误差可定义为
\begin{equation}
  e_{\mathrm{gc}}(t)=\dist\bigl(p_c,\ell(t)\bigr).
  \label{eq:ch1_gc_error}
\end{equation}
对于微创手术，$p_c$ 可表示穿刺点；对于工业场景，$p_c$ 可表示导向孔、夹具中心、深腔入口或固定通道中心。式\eqref{eq:ch1_gc_error} 说明，本文讨论的几何约束并不局限于医学远心运动约束，而是可覆盖一类固定点--工具轴线约束。机器人运动学、动力学和雅可比转置映射为这类任务提供了建模基础\cite{siciliano2016handbook,khatib1987operational}。

为了更清楚地说明本文研究对象的实际来源，表\ref{tab:ch1_application_scenarios} 给出了几类典型约束接触任务。可以看到，不同场景的工具、工件和评价指标并不完全相同，但控制层面的共性非常明显：一方面，机器人需要沿给定路径推进，保证作业覆盖、加工一致性或操作精度；另一方面，末端接触力必须被限制在合理范围内。更重要的是，操作者往往不是简单地“开关”机器人，而是在任务执行过程中根据视觉观察、经验判断或工艺要求对轨迹进行局部修正。若控制器只关注自主轨迹，人的修正难以进入任务；若控制器完全服从人，接触力又可能失去边界保护。

\begin{table}[H]
  \centering
  \caption{典型工业约束接触任务及其控制矛盾}
  \label{tab:ch1_application_scenarios}
  \begin{tabular}{p{0.20\textwidth}p{0.25\textwidth}p{0.25\textwidth}p{0.20\textwidth}}
    \toprule
    应用任务 & 操作者作用 & 机器人作用 & 主要风险 \\
    \midrule
    柔性密封条检测 & 根据外观缺陷和检测信号修正扫描区域 & 保持稳定轻接触并沿表面扫描 & 接触丢失或压伤软质材料 \\
    复合材料壳体打磨 & 根据局部余量调整加工路径 & 保持恒力打磨和重复覆盖 & 过磨、欠磨或工具冲击 \\
    曲面抛光与去毛刺 & 选择重点区域和局部边缘 & 维持切向轨迹和法向接触力 & 表面质量不均、力矩峰值过大 \\
    受限通道操作或手术训练 & 提供目标修正和安全监督 & 保持工具通道约束和末端稳定 & 几何约束超限、组织/工件损伤 \\
    \bottomrule
  \end{tabular}
\end{table}

表\ref{tab:ch1_application_scenarios} 也解释了为什么本文不采用单一场景叙事。微创手术训练任务便于构造明确的固定点--轴线几何约束，柔性工件打磨和表面检测则更能体现工业落地需求。二者在数学建模上都需要处理参考轨迹、人类修正、接触力和约束误差。因此，本文把具体实验平台看作验证载体，把“约束接触任务中的动态仲裁控制”作为更一般的研究对象。

进一步从控制系统结构来看，本文所关注的任务至少包含四类信息流。第一类是机器人自身状态，包括关节位置、速度、末端位姿和雅可比矩阵，这些信息决定了机器人是否能够在当前构型下执行期望运动。第二类是操作者输入，包括主端位移、速度、交互力或手柄命令，它们反映了人的即时修正意图。第三类是环境交互信息，包括接触力、接触法向、表观刚度和局部压入深度，它们反映了任务是否仍处于安全接触状态。第四类是任务约束信息，包括工件边界、禁入区域、固定通道、工具长度和安全力范围。传统单回路控制器通常只处理其中一两类信息，而本文希望构造的是能同时接收并解释这些信息的分层控制框架。

这种分层并不是人为复杂化，而是由任务本身决定的。操作者输入通常发生在较低频率，其作用是改变任务目标或局部路径；接触力反馈则需要在更高频率下进入执行层，用于抑制过压和接触丢失；几何约束映射位于二者之间，它把工具尖端期望转换为机器人实际执行的法兰或关节命令。如果把人的输入直接加到力矩层，容易破坏接触稳定性；如果把力安全完全放在参考层，系统又可能响应过慢。因此，本文后续章节采用“参考层仲裁--执行层仲裁--约束映射”的结构：人机意图主要决定期望轨迹，接触安全主要决定力--位优先级，几何约束决定执行映射。

以柔性密封条检测为例，机器人自主轨迹可以由 CAD 模型或示教路径生成，但真实工件在装夹后往往存在毫米级偏差。操作者看到检测探头偏离边缘时，会给出横向修正；这类切向修正应被控制器接受，否则检测区域不完整。与此同时，若操作者沿法向误推，使探头进一步压入软质材料，控制器就必须限制该输入，因为过大的法向力会改变检测信号甚至损伤工件。类似地，在复合材料打磨中，操作者可能要求机器人绕开局部缺陷或增加某一区域覆盖次数，但机器人仍需保持恒定法向力，避免局部过磨。可以看出，人类意图和接触安全并非同一层面的变量，它们需要被分别建模。

上述分析也说明，本文研究并不追求完全自主替代人类。对于很多高价值、小批量或安全敏感任务，完全自主系统的部署成本和风险仍然较高。更现实的形态是，机器人承担可重复、高精度和高带宽的执行部分，人类负责目标修正、异常判断和策略监督。这样的系统需要回答三个问题：人什么时候应该影响轨迹？人影响后的轨迹是否安全？机器人在执行该轨迹时应更重视位置还是力？第三章、第四章和第五章分别对应这三个问题的递进回答。

从实验评价角度看，这类系统也不能只使用单一指标判断好坏。若只比较轨迹误差，机器人可能通过增大位置刚度获得较小误差，但接触力随之增大；若只比较力误差，机器人可能停留在局部区域而不完成扫描路径；若只比较操作者输入响应速度，系统又可能过度服从人类瞬时输入。因此，本文后续章节同时关注轨迹精度、力安全、几何约束、输入平滑性和实时性。轨迹精度体现任务完成质量，力安全体现工件和组织保护，几何约束体现工具通道合规，输入平滑性体现人机协同体验，实时性则决定方法能否真正进入机器人控制循环。这样的评价体系更接近真实工业和手术训练场景。

在控制器设计上，本文也尽量避免把所有复杂性推给参数调节。例如，若把人机仲裁、力--位仲裁和安全边界全部压缩为一个总权重，确实可以写出更短的公式，但该权重的物理含义会变得模糊。调试时很难判断是人的意图识别出了问题，还是接触力安全判断出了问题，或者是底层力--位控制器增益不合适。本文采用多个变量分层描述：$\alpha_{\mathrm{HR}}$ 表示人机监督权重，$\beta$ 表示参考层直接遥操作参考比例，$\alpha_{\mathrm{FP}}$ 表示力--位优先级，$\rho_F$ 表示接触力安全裕度。这样虽然变量数量略多，但每个变量都有清楚的数学定义和物理解释，便于理论分析和实验诊断。

还需要指出，工业落地中的“安全”并不只意味着不发生碰撞。对于接触式打磨和检测任务，轻微力越界也可能造成表面质量下降；对于手术训练，短时力冲击可能让受训者形成错误操作习惯；对于力传感器和细长工具，过大接触力还可能导致传感器饱和、工具弯曲或驱动线缆受力异常。因此，本文所说的接触安全不是事后急停，而是希望控制器在力接近边界之前就改变行为趋势。第四章和第五章中的力安全裕度仲裁正是服务于这一目的：它把“离边界还有多远”转化为连续控制权重，而不是等到越界后再切换模式。




综上，本文研究具有三方面意义。第一，从理论上，本文将人机共享控制、力--位协同控制和几何约束执行放入合作博弈框架中，避免将多目标控制退化为经验权重调节。第二，从方法上，本文区分人机意图仲裁和力--位安全仲裁：前者作用于参考轨迹层，后者作用于执行控制层，二者物理含义清晰。第三，从应用上，本文面向工业柔性接触任务和机器人辅助手术训练任务，强调方法的可解释性、实时性和可落地性。与端到端学习方法相比，本文更关注在有限传感器、有限模型精度和安全边界明确的条件下，如何构造可分析、可调试、可验证的控制器。

\section{国内外研究现状}
\label{sec:ch1_related_work}

\subsection{人机共享控制与意图理解}
\label{subsec:ch1_shared_autonomy}

共享控制的基本思想是，人和机器人共同影响系统行为。当机器人能够可靠判断操作者目标时，可以提高自主程度以降低人的操作负担；当目标不确定或风险较高时，应降低自主介入，避免错误接管。早期共享控制通常依赖固定候选目标、固定权重或启发式规则。随着任务复杂度提高，研究重点逐渐转向长期交互、个性化辅助、认知负荷、信任校准和上下文感知通信等问题\cite{matsumoto2022shared,li2023trust,pan2024effects,luo2024customizable}。从控制权分配角度看，连续角色适应方法说明，人和机器人之间不必只在“人主导”和“机器人主导”之间离散切换，而可以依据交互力和任务状态在线平滑移动角色边界\cite{li2015continuousrole}；维度相关共享自主进一步指出，单一标量权重有时会掩盖不同自由度上的人机分歧，因此仲裁变量最好和任务约束、运动方向以及操作者接受程度联系起来\cite{bowman2023dimensionspecific}。

Jonnavittula 和 Losey 等提出的重复交互共享自主方法强调，机器人不应总是假设候选目标已知，而应从用户反复执行的任务中学习辅助策略\cite{jonnavittula2021learning,jonnavittula2022sari}。这类研究说明，人类意图不是一次性识别问题，而是长期交互过程中的持续估计问题。Yousefi 等进一步从人类在环策略微调和共享自主对齐角度讨论复杂任务中的多层交互，说明不同任务层级的自主辅助需要不同抽象程度的参考生成和策略更新机制\cite{yousefi2025sharedautonomy}。Pan 等通过用户实验分析共享控制中的认知负荷和信任关系，提示共享控制不能只优化轨迹误差，还需要考虑操作者对系统行为的理解和接受程度\cite{pan2024effects}。

人类输入的困难在于其含义不唯一。短时、小幅输入可能是误触、震颤，也可能是精细修正；持续、大幅输入可能表示明确接管或目标重定向；输入方向与机器人自主方向一致时，可能只是希望加快任务进度，而相反方向输入可能代表对自主目标的否定。轨迹变形方法将物理交互解释为对未来轨迹的局部修正\cite{losey2018trajectory}，概率意图识别方法则通过候选目标后验概率描述操作者意图\cite{jain2020probabilistic}。进一步地，Losey 等将物理交互看作人的目标信息，而不只是外部扰动，强调机器人应当从人的纠正中更新后续任务行为\cite{losey2022physicalcommunication}。这些方法为本文第三章的人机意图特征提取提供了基础：操作者输入应当先被解释为监督意图，再进入共享参考和博弈代价，而不是直接作为底层控制命令。

从实际系统实现来看，共享控制还需要处理三个细节。第一，人的输入通道通常带有低频意图和高频扰动的混合成分。对于主从设备，手部微小震颤、设备摩擦和视觉反馈延迟都会造成输入波动；对于工业手柄或力引导设备，操作者可能在观察工件表面时产生短时试探输入。第二，人类输入具有方向性。沿工件表面切向的输入往往表示希望改变扫描区域，而沿法向的输入更可能改变压入深度和接触力。第三，共享控制的输出必须平滑。触觉共享控制研究较早就强调，控制权若能通过力反馈连续变化，操作者更容易理解自动化边界并保持在环参与\cite{abbink2011hapticshared}；面向手术操作的活动感知共享控制和 STAR 置信共享控制也说明，权重切换需要依赖任务状态、控制置信度或活动识别，而不是只依赖一个瞬时输入量\cite{su2022activityaware,saeidi2018starshared}。因此，本文第三章和第五章均采用“特征提取--模糊/规则仲裁--滤波限幅”的流程，而不是把瞬时输入幅值直接映射为控制权重。


在安全共享自主方面，He 等提出的 Barrier Pair 方法直接关注未知人类输入给共享控制带来的安全挑战，表明共享自主系统必须在辅助人类和满足安全约束之间进行平衡\cite{he2021barrier}。控制屏障函数和安全学习相关研究也说明，安全约束不能完全依赖事后限幅，而应在控制器设计时进入约束、证书或代价结构；面向动态风险区域的人机协作最优运动规划研究也从任务规划层面验证了风险感知约束的重要性\cite{ames2019cbf,xiao2019hocbf,dawson2021rnlbf,cohen2021safeexploration,brunke2022safelearning,li2025optimal}。这些工作共同提示本文：在人机协同接触任务中，人类输入应当经过安全过滤和连续仲裁，不能直接越过接触力边界和几何约束。

\subsection{微分博弈与动态仲裁控制}
\label{subsec:ch1_game_control}

人机协同控制天然具有多目标属性。人侧目标、机器人自主目标、位置跟踪目标、接触力目标和安全约束往往同时存在，并且不能简单相加。微分博弈和最优控制为这种多参与方、多目标问题提供了系统化语言。经典动态博弈理论给出了参与方策略、代价函数和均衡解的数学定义\cite{basar1999dynamic,engwerda2005lqdg}；线性二次最优控制理论则为代数黎卡提（Riccati）方程、状态反馈和稳定性分析提供了基础工具\cite{anderson2007optimal}。

近年来，微分博弈逐渐进入物理人机交互研究。Li 等将微分博弈用于多样化物理人机交互任务，展示了博弈建模在理解人和机器人作用力分配方面的潜力\cite{li2019differential}。Franceschi 等进一步研究了合作微分博弈在人机协作到达任务中的应用，并比较了合作博弈、非合作博弈和 LQR 控制的差异；其后续工作将微分博弈用于人机角色仲裁和自适应阻抗控制，说明博弈权重可以被解释为人和机器人在任务目标中的相对作用\cite{franceschi2022adaptive,franceschi2023role,franceschi2023modeling}。此外，人类运动意图预测和自适应阻抗控制研究也强调，在物理交互中，控制器需要同时估计人的目标和调节机器人柔顺性\cite{franceschi2023predicting,ferrante2022adaptive}。

与启发式加权不同，合作博弈控制的优势在于，权重不是孤立的经验参数，而是在帕累托（Pareto）前沿上选择工作点的变量。帕累托前沿表示这样一组折中解：若希望继续改善某一目标，就必须牺牲至少另一个目标。对于第三章，$\alpha_{\mathrm{HR}}$ 描述人侧目标与机器人自主目标之间的相对权重；对于第四章，$\alpha_{\mathrm{FP}}$ 描述位置目标与力目标之间的相对优先级。二者都可以通过加权代价和黎卡提方程得到反馈律，但二者物理含义不同，不应混用。本文第五章进一步强调这种分层思想：人机仲裁影响期望轨迹，力--位仲裁影响执行层反馈增益。

在具体推导中，本文主要采用线性二次型或局部线性化后的二次型代价。这一选择并不是因为真实任务完全线性，而是因为它具有三个工程优势。首先，二次代价可以直接表达位置误差、速度误差、力误差和控制输入能量等物理量，参数含义清楚，便于调试。其次，在线性化模型下，黎卡提方程能够给出解析结构明确的状态反馈律，不需要在每个控制周期求解高维非线性规划。再次，若环境参数或仲裁权重变化较慢，可以通过离线增益表和在线插值实现实时控制。对于百赫兹量级的机器人力矩控制或阻抗控制，这一点非常重要。换句话说，本文选择合作博弈并非只追求形式上的统一，而是希望在可解释建模、实时求解和稳定性分析之间取得平衡。


博弈理论还与学习和教学任务相结合。Yu 等将机器人教学形式化为合作 Markov 博弈，说明机器人可以在交互过程中承担教学或辅助角色\cite{yu2023coach}。不过，这类学习式或教学式方法通常更关注任务策略层，而本文关注的是低层连续控制中的轨迹、力和约束协同。因此，本文采用合作博弈和最优控制相结合的形式，力求在理论可分析性和工程实时性之间取得平衡。

\subsection{约束接触操作与力--位协同控制}
\label{subsec:ch1_force_position}

接触控制是机器人从自由空间运动走向真实操作任务的关键。混合位置/力控制通过选择矩阵区分运动控制方向和力控制方向，阻抗控制通过期望惯性、阻尼和刚度描述机器人与环境之间的柔顺关系\cite{raibert1981hybrid,hogan1985impedance}。这些方法解释了一个基本事实：在接触任务中，位置和力不是独立变量。沿自由方向，机器人应尽量跟踪轨迹；沿约束方向，机器人应调节接触力或表现出合适的柔顺性。在人机协同或手术辅助场景中，虚拟夹具和主动约束常被用来把“不能进入的区域”“应该沿着的方向”转化为控制或触觉反馈边界\cite{bowyer2014virtualfixtures,marinho2019dynamicactive}，这与本文将安全边界写入仲裁变量的思想是一致的。

近年来，接触丰富操作研究开始从纯位置轨迹学习转向显式力信息利用。Beltran-Hernandez 等研究了位置控制机器人上的学习式力控制，将传统力控结构与强化学习结合，用于真实接触任务\cite{beltran2020learning}。ForceMimic 从人类示教中采集力--运动信息，并通过混合力--位模仿学习增强接触任务鲁棒性\cite{liu2024forcemimic}。FORGE 将最大允许力阈值引入不确定接触操作策略学习，用于避免过于激进的接触行为\cite{noseworthy2024forge}。腿足机器人操作、多接触全身力控和柔性力传感抓取垫等研究也说明，力调节正在成为通用机器人操作中的重要组成部分\cite{portela2024learning,rouxel2023seiko,han2024grippingpads}。

同时，大模型和扩散策略等方法显著提升了机器人对复杂任务的感知和策略生成能力。Diffusion Policy 在多类机器人操作基准中展示了扩散模型在动作序列生成中的优势\cite{chi2023diffusion}；RT-2 和 PaLM-E 等视觉--语言--动作或具身多模态模型说明，机器人可以从大规模视觉语言数据中获得更强的语义泛化能力\cite{brohan2023rt2,driess2023palme}。然而，这些方法并不自动解决接触力安全和稳定性问题。对于打磨、检测、插入和软组织操作这类安全边界明确的任务，仅依靠端到端策略输出仍然难以解释“为什么此时应该退让”“为什么人的法向输入应被抑制”。因此，本文仍采用模型控制和博弈控制作为底层核心，同时吸收学习方法中关于人类意图、任务上下文和力信息重要性的认识。

对于工业约束接触任务，另一个容易被忽略的问题是“切向任务”和“法向安全”的差异。以表面扫描为例，切向方向决定检测覆盖率和路径效率，法向方向决定压入深度和接触力；以抛光为例，切向速度影响覆盖均匀性，法向力影响材料去除率和表面质量。如果使用一个笛卡尔位置误差统一驱动所有方向，当接触力接近上界时，控制器可能为了跟踪空间轨迹继续压入；如果使用纯力控，又可能牺牲切向轨迹覆盖。因此，本文第五章在第四章力--位协同控制的基础上进一步引入切向/法向分解：人的安全切向修正可以被接受，人的危险法向压入则应被抑制。这种设计更贴近工业现场的操作逻辑。


工业表面处理任务进一步体现了力--位协同的必要性。RoboGrind 面向风电叶片修复等表面处理任务，结合感知、交互式程序生成和力控执行，说明实际工业场景不只是轨迹跟踪，还包括缺陷识别、区域选择和力控加工\cite{alt2024robogrind}。增强现实示教编程研究则指出，小批量和高变化产品中的机器人部署仍受编程复杂度限制，非专家需要通过直观交互影响机器人轨迹\cite{dogangun2025rampa}。安全且高效的人机协作研究表明，协作空间和安全边界若设置过于保守，会明显降低生产率；若过于激进，又会带来人机协作风险\cite{kamezaki2024dynamicworkspace}。这些工作都从应用角度支持本文第五章的设计：工业约束接触任务需要同时考虑人的在线轨迹干预和接触安全边界。类似逻辑在手术力反馈任务中也很常见，例如经皮肾镜训练和针抓取优化中，触觉提示既降低操作者负担，又保留了人的最终决策权\cite{wilz2021pcnlhaptic,selvaggio2019hapticneedle}。

\subsection{机器人辅助手术与安全自主操作}
\label{subsec:ch1_surgical_industrial}

机器人辅助手术是约束接触控制的重要应用方向。微创手术系统通过长工具进入人体，工具运动受到穿刺孔和器械结构约束；手术组织具有柔性和变形特性；医生需要保留监督和接管能力。早期关于手术机器人自主性的综述已经指出，自主手术需要在感知、规划、控制、力反馈和临床安全之间建立闭环\cite{yip2017autonomy,simaan2018medical}。近五年，软组织自主缝合和吻合研究表明，机器人可以在特定结构化任务中达到较高的操作一致性，但也暴露出软组织变形建模、视觉跟踪、力感知和安全验证方面的挑战\cite{saeidi2022star}。在协作模式上，多边手术操作框架将遥操作、共享控制、监督控制、交换控制和全自主放在同一任务中比较，说明手术场景中的人机协作不是单一模式选择，而是不同任务阶段、不同风险水平下的控制权组织问题\cite{nichols2016multilateral}。

这类挑战与本文所考虑的远心运动或固定通道约束密切相关。RCM 约束最早主要服务于微创手术中的穿刺孔保护，但近年来也被总结为一类围绕远端固定点运动的通用机器人约束\cite{zhang2024rcmreview}。长工具绕固定点转动时，工具尖端位置、法兰位置和工具姿态并不是独立变量。为了到达某个末端位置，机器人可能需要改变法兰姿态；姿态变化又会改变工具与组织或工件之间的接触角度，进而改变法向力。现有 RCM 控制研究从关节空间导纳模型、动力学扭矩控制和被动导纳约束等角度给出了不同实现方式\cite{kastritsi2021rcmcontroller,minelli2021dynamicrcm,kastritsi2022passiveadmittance}。因此，如果仅在工具尖端空间设计力控，再简单通过杠杆比例映射到法兰空间，容易产生物理含义不清的问题。本文第三章和第四章均采用工具期望到法兰期望的正向几何映射思想，第五章则将这一思想推广到工业导向孔、固定通道和约束接触任务中。


在安全辅助方面，增强现实和情境感知力控制为手术机器人提供了两类思路。增强现实框架通过重建和配准关键解剖结构，帮助操作者避免器械与血管等高风险区域碰撞\cite{chen2023markerless}；情境感知力控制和触觉辅助框架则通过力控和虚拟夹具降低钻削等任务中的不期望接触力\cite{ishida2024force,ishida2024haptic}。在受限内腔操作中，共享自主柔性机械臂和儿童内镜缝合虚拟夹具研究进一步说明，医生可保留路径定义或监督角色，而机器人负责约束环境中的局部轨迹规划、碰撞规避和精细执行\cite{ma2019endoluminalshared,marinho2020virtualfixture}。与此同时，面向手术自主性的基础模型研究指出，虽然多模态大模型有望增强手术机器人的场景理解和任务规划能力，但数据稀缺、软组织仿真差距和临床安全要求仍然限制其直接落地\cite{schmidgall2024foundation,zargarzadeh2024bloodsuction}。

上述手术机器人研究与工业约束接触任务具有共性。两者都需要在通道约束、接触力、人的监督和机器人自主之间取得平衡。不同之处在于，工业任务更强调产线效率、表面质量和可重复部署；手术任务更强调组织安全、医生监督和极端风险避免。因此，本文没有将方法限制在单一手术场景，而是采用更一般的“工业约束接触任务”作为第五章的统一落地背景。这样既可以支撑前几章在手术训练平台上的验证，也能说明方法在柔性工件检测、打磨和接触式加工中的推广潜力。


综上，现有研究已经分别在共享自主、博弈控制、力控安全和手术机器人等方向取得了重要进展，但对本文关注的问题仍存在不足。表\ref{tab:ch1_gap_summary} 对这些不足和本文对应研究内容进行了概括。该表也给出了全文第三、四、五章之间的逻辑关系：第三章解决人如何影响机器人目标，第四章解决接触力如何影响力--位优先级，第五章则进一步解决二者同时存在时如何分层协调。

从文献发展趋势看，学习方法和大模型方法正在快速提升机器人对复杂场景的理解能力，安全控制方法正在提升机器人满足约束的理论能力，力感知操作方法正在提升机器人处理接触任务的鲁棒性。但是，这些方向之间仍然存在接口问题。学习方法生成的轨迹需要经过安全约束检查，安全控制方法需要明确哪些输入代表人的有效意图，力控方法需要知道何时应牺牲部分位置精度以换取接触安全。本文的工作可以理解为对这些接口问题的控制层回答：不直接讨论如何训练大模型或如何构建完整工业感知系统，而是在给定人类输入、接触力反馈和机器人状态的条件下，建立一个可解释、可计算、可稳定性分析的仲裁控制框架。

因此，本文创新性的重点并不在于提出一个全新的机器人硬件平台，也不在于用一个端到端模型覆盖所有任务，而在于明确区分不同控制变量的物理含义。$\alpha_{\mathrm{HR}}$ 表示人机监督权重，它回答“期望轨迹更接近人还是机器人”；$\alpha_{\mathrm{FP}}$ 表示力--位优先级，它回答“执行控制更重视位置还是接触力”；$\rho_F$ 表示接触力安全裕度，它回答“当前力是否仍处在安全区间中部”。这些变量分工明确后，控制器设计、参数调试和实验分析都会更加清晰。


\begin{table}[H]
  \centering
  \caption{相关研究不足与本文研究内容的对应关系}
  \label{tab:ch1_gap_summary}
  \begin{tabular}{p{0.25\textwidth}p{0.34\textwidth}p{0.30\textwidth}}
    \toprule
    研究方向 & 主要不足 & 本文对应内容 \\
    \midrule
    共享自主与意图识别 & 多关注目标推断和轨迹辅助，对接触力安全边界处理不足 & 第三章建立人机意图仲裁，第五章加入安全参考投影 \\
    微分博弈与动态仲裁 & 常将不同物理意义的权重统一处理，容易造成解释混淆 & 区分 $\alpha_{\mathrm{HR}}$ 与 $\alpha_{\mathrm{FP}}$ 的作用层级 \\
    接触力控制与学习操作 & 强调力控或策略学习，但较少考虑人类在线修正期望轨迹 & 第四章构造力--位协同控制，第五章加入人影响参考 \\
    手术与工业落地应用 & 场景约束真实复杂，但底层控制常缺少统一可推演模型 & 以约束接触任务统一几何、力和人机监督三类约束 \\
    \bottomrule
  \end{tabular}
\end{table}

\section{本文研究内容与创新点}
\label{sec:ch1_contributions}

综合上述研究现状可以看到，现有方法仍存在三个明显不足。第一，许多共享控制方法关注人类意图识别和目标预测，但对接触力安全边界考虑不足。第二，许多力控或接触丰富操作方法关注力--位协同，但没有充分处理人类在线修改期望轨迹的问题。第三，部分博弈控制方法能够描述多目标权衡，但容易将不同物理含义的权重合并为单一参数，导致模型解释和工程实现都不够清晰。本文围绕这三个问题，构建从刚性约束共享控制、柔性接触力--位协同到工业约束接触双仲裁控制的递进式研究框架。

第一，本文研究面向刚性几何约束操作的人机意图驱动合作博弈共享控制方法。针对长工具通过固定通道或远心运动点执行任务时，操作者输入中同时包含短时扰动、局部修正和目标重定向的问题，本文提取主端交互强度、干预持续时间和空间风险等特征，构造人机监督权重 $\alpha_{\mathrm{HR}}$。在控制层面，将人侧监督目标与机器人自主目标写入合作博弈代价函数，并通过几何映射将任务空间控制输入转换为满足工具通道约束的法兰执行命令。该部分的创新点在于：将操作者输入解释为对共享目标的监督调节，而不是直接等同于从端运动命令，从而提高约束操作中的稳定性和可解释性。

第二，本文研究面向柔性接触任务的折扣帕累托力--位协同控制方法。针对柔性环境中位置误差和接触力误差强耦合的问题，本文构造包含位置误差、速度误差、接触力误差和力误差泄漏积分的增广状态 $\xi$，将位置目标和力目标作为两个协同参与方，在折扣型帕累托代价下求解反馈控制律。与直接使用环境刚度作为安全风险指标不同，本文引入接触力安全裕度 $\rho_F$ 和边界方向变量 $s_F$，根据实际交互力与上下界的距离调节力--位优先级 $\alpha_{\mathrm{FP}}$。该部分的创新点在于：将接触力安全边界显式嵌入仲裁机制，使控制器能够在轨迹推进和力安全之间动态调整。

第三，本文研究面向工业约束接触任务的双仲裁合作博弈控制方法。实际工业任务中，操作者既可能希望改变机器人扫描轨迹，又可能无意中引入法向过压风险。为此，本文提出参考层人机仲裁和执行层力--位仲裁相分离的框架。$\alpha_{\mathrm{HR}}$ 只作用于人影响后的安全期望轨迹生成，$\alpha_{\mathrm{FP}}$ 只作用于底层力--位反馈增益选择。这样既继承了第三章的人机意图理解能力，又继承了第四章的接触安全调节能力，同时避免构造过高维的黎卡提增益表。该部分的创新点在于：把“人可修正轨迹”和“系统守住接触力边界”统一到一个可实现的分层博弈控制框架中。

本文的主要创新可概括为以下三点：
\begin{enumerate}[label=(\arabic*)]
  \item 提出面向刚性约束操作的人机意图动态仲裁合作博弈方法，区分人侧监督参考、机器人自主参考和人机博弈权重，使主端输入能够以连续、平滑和可解释方式影响共享控制目标。
  \item 提出面向柔性接触任务的交互力安全裕度仲裁方法，将接触力到上下安全边界的距离引入力--位帕累托控制器，使控制器能够根据实际接触风险在位置跟踪和力调节之间切换。
  \item 提出面向工业约束接触任务的双仲裁分层控制框架，将 $\alpha_{\mathrm{HR}}$ 限定在参考轨迹生成层，将 $\alpha_{\mathrm{FP}}$ 限定在执行控制层，从而在不显著增加在线求解复杂度的情况下实现人机意图、接触安全和几何约束的统一协调。
\end{enumerate}

需要说明的是，本文并不试图直接解决所有机器人自主操作问题，也不假设机器人可以完全准确预测人的所有意图。本文关注的是一种更接近当前工程落地阶段的协同控制形态：机器人承担稳定、重复和高带宽控制任务，操作者提供监督、修正和异常判断，控制器通过动态仲裁在人的输入、机器人的自主目标和接触安全边界之间进行协调。这一定位与工业柔性接触任务和机器人辅助手术训练平台的实际需求相符合。

\section{全文组织结构}
\label{sec:ch1_organization}

全文共分为六章，各章内容安排如下。

第一章为绪论。本章首先介绍人机协同机器人、工业约束接触任务和机器人辅助手术中的应用背景，分析操作者意图、几何约束和接触力安全之间的矛盾；随后从共享控制、微分博弈、力--位协同控制和落地应用四个方面综述国内外研究现状；最后给出本文研究内容、创新点和章节安排。

第二章为预备知识。该章介绍机器人运动学与动力学、约束运动与接触建模、最优控制与微分博弈、稳定性与鲁棒性分析等基础内容，并统一全文使用的主要符号，为后续章节的建模和理论分析奠定基础。

第三章研究面向刚性几何约束操作场景的意图驱动合作博弈共享控制方法。该章从刚性约束任务出发，建立人机监督共享控制问题，设计主端意图特征提取和动态仲裁机制，并在合作博弈框架下求解共享控制律，最后通过几何约束执行映射实现工具到法兰的控制命令转换。

第四章研究面向柔性接触任务的折扣帕累托力--位协同控制方法。该章建立含残差项的柔性接触增广系统，构造位置目标和力目标的折扣帕累托代价函数，推导相应反馈控制律，并设计基于交互力安全裕度的力--位优先级仲裁机制。

第五章研究面向工业约束接触任务的双仲裁合作博弈控制方法。该章将第三章的人机意图仲裁与第四章的力--位安全仲裁统一到分层框架中，提出参考层安全轨迹融合、执行层力--位优先级调节以及相应仿真验证方案，用于展示方法在通用工业柔性接触任务中的实用性。

第六章总结全文工作，并讨论未来可进一步开展的研究方向，包括真实力传感器实验、多模态意图识别、复杂曲面工件验证、严格控制屏障函数安全约束以及面向真实手术器械的力感知与控制扩展。
